\section*{2 Preliminaries}
\subsection*{2.1 Definitions and notation}
\paragraph{Simultaneous messages protocols}
A (two-party) simultaneous messages (SM) protocol consists of two players, which we refer to as Alice and Bob, and a referee, which we refer to as Carol, that together compute a function. Alice and Bob each send a message, which is based on the input, to Carol, who then computes a function of the two messages received. Formally, we have the following definitions:
\setcounter{nttheorem}{2}
\begin{ntdefinition}[Author Definition3: Simultaneous messages protocols]
Let $X,Y,\widehat X,\widehat Y,Z$ be finite nonempty sets. A simultaneous messages protocol (shortly, SM protocol or SMP) $P$ is a triplet of functions $(A,B,C)$, where $A:X\to\widehat X$, $B:Y\to\widehat Y$, and $C:\widehat X\times\widehat Y\to Z$. We call $C$ the referee function.
\end{ntdefinition}
\begin{ntdefinition}[Author Definition4: SM protocol admittance]
Let $f:X\times Y\to Z$ be a function. We say that $f$ admits an SM protocol $(A,B,C)$ if $f(x,y)=C(A(x),B(y))$ for every $(x,y)\in X\times Y$. In that case, we also say that $(A,B,C)$ computes $f$.
\end{ntdefinition}
\paragraph{Inner product space of Boolean functions}
For the purpose of utilizing Fourier analysis, we will consider the inner product space of all functions $\{-1,1\}^n\to\mathbb R$ with the following inner product:
\[\langle f,g\rangle=\mathbb E_{x\sim\{-1,1\}^n}[f(x)g(x)]=\frac1{2^n}\sum_{x\in\{-1,1\}^n}f(x)\cdot g(x).\]
It is a known fact that the aforementioned inner product space has as orthonormal basis the set of all parity functions $\{\chi_S\}_{S\subseteq[n]}$, defined by $\chi_S(x)=\prod_{i\in S}x_i$.
\paragraph{The Inner Product function}
The inner product modulo2 function $\operatorname{IP}:\{0,1\}^n\times\{0,1\}^n\to\{0,1\}$ is defined by
\[\operatorname{IP}(x,y)=\sum_{i=1}^nx_iy_i\pmod2.\]
Note that each $x$, when fixed, corresponds to a parity function on a subset $S$ of $y$'s coordinates, a subset which is determined by $x$. This correspondence is actually a bijection mapping elements of $\{0,1\}^n$ to subsets of $2^{[n]}$. Thus, when we switch to $\pm1$ notation, we can identify each $x\in\{-1,1\}^n$ with the parity function $\chi_{S(x)}(y)$ that results when fixing $x$ in $\operatorname{IP}(x,y)$.
\paragraph{One-to-one condition}
Some functions have a certain property, formally defined below, rendering them harder to compute with preprocessing. Restricting attention to these functions seems to help in obtaining lower bounds.
\begin{ntdefinition}[Author Definition5: One-to-one condition]
Let $f:X\times Y\to Z$ be a function. We say that $f$ satisfies the left one-to-one condition if for every $x\ne x'\in X$ there exists $y\in Y$ such that $f(x,y)\ne f(x',y)$. Similarly, we say that $f$ satisfies the right one-to-one condition if for every $y\ne y'\in Y$ there exists $x\in X$ such that $f(x,y)\ne f(x,y')$. Finally, we say that $f$ satisfies the one-to-one condition if $f$ satisfies both the left and right one-to-one conditions.
\end{ntdefinition}
\begin{ntproposition}[Author Proposition6: One-to-one preprocessing]
Let $f:X\times Y\to Z$ be a function. Then:
\begin{itemize}
\item If $f$ satisfies the left one-to-one condition, then for every SM protocol $(A,B,C)$ that $f$ admits, $A$ computes a one-to-one mapping.
\item If $f$ satisfies the right one-to-one condition, then for every SM protocol $(A,B,C)$ that $f$ admits, $B$ computes a one-to-one mapping.
\end{itemize}
\end{ntproposition}
\subsection*{2.2 Known facts}
The following are known facts we will need later.
\begin{nttheorem}[Author Theorem7: Tal's LMN improvement (LMNT), {\cite{ref24,ref30}}]
Let $f$ be a Boolean function with $n$ variables computable by an unbounded fan-in circuit of depth $h$ and size $M$, and let $t$ be any integer. Then,
\[\|f^{\geq t}\|^2\leq2\cdot2^{-t/O_h(\log M)^{h-1}}.\]
\end{nttheorem}
\begin{ntlemma}[Author Lemma8, Lemma3.6 in {\cite{ref21}}]
Let $H:[0,1]\to\mathbb R$ be the binary entropy function defined by $H(p)=-p\log p-(1-p)\log(1-p)$. Then, for any $0<a\leq1/2$ and $n\in\mathbb N$,
\[\binom n{\leq an}\leq2^{H(a)n}.\]
\end{ntlemma}
